\documentclass[aspectratio=169,11pt]{beamer}

\usetheme{Madrid}
\usecolortheme{seahorse}

\usepackage{fontspec}
\usepackage{booktabs}
\usepackage{array}
\usepackage{hyperref}

\setsansfont{Latin Modern Sans}
\setmonofont{Latin Modern Mono}
\hypersetup{colorlinks=true, urlcolor=blue}
\setbeamertemplate{navigation symbols}{}
\setbeamertemplate{footline}[frame number]

\title{BDS Projekat 1}
\subtitle{AIS analitika nad HDFS + Spark}
\author{Mihajlo Madic 2119}
\institute{Big Data Systems}
\date{}

\newcommand{\code}[1]{\texttt{\small #1}}

\begin{document}

\begin{frame}
  \titlepage
\end{frame}

\begin{frame}{Dataset: Piraeus AIS}
  \begin{itemize}
    \item Korišćen je javni dataset \textit{The Piraeus AIS Dataset for Large-scale Maritime Data Analytics}, objavljen na Zenodo repozitorijumu.
    \item Dataset sadrži AIS poruke brodova iz šire zone luke Pirej u Grčkoj.
    \item AIS zapisi opisuju kretanje plovila kroz atribute kao što su vreme, identifikator plovila, geografska pozicija, brzina, kurs i pravac.
    \item Originalni dataset je dovoljno veliki za Big Data scenario; za ovaj prototip korišćen je kontrolisani podskup radi bržeg razvoja i ponovljivog benchmark-a.
    \item Izvor: \url{https://zenodo.org/records/5792100}
  \end{itemize}
\end{frame}

\begin{frame}{Zadatak}
  Implementirati ceo sistem:
  \begin{enumerate}
    \item priprema podataka i upload na HDFS
    \item Spark aplikacija sa CLI upitima (\code{count}, \code{show}, \code{stats})
    \item benchmark poređenje režima izvršavanja
  \end{enumerate}

  \vspace{0.4cm}
  Rezultati projekta:
  \begin{itemize}
    \item reproducibilna infrastruktura kroz Docker Compose
    \item merenja performansi
    \item izveštaji sa rezultatima
  \end{itemize}
\end{frame}

\begin{frame}{Obim i podaci}
  Korišćen podskup:
  \begin{itemize}
    \item raw ulaz: \code{\textasciitilde 971.54 MiB} (jul + decembar 2018 CSV)
    \item preprocesirano: \code{\textasciitilde 191.60 MiB} (Parquet)
  \end{itemize}

  \vspace{0.35cm}
  Relevantne HDFS putanje:
  \begin{itemize}
    \item \code{/bds/proj1/raw}
    \item \code{/bds/proj1/processed/clean}
    \item \code{/bds/proj1/processed/quarantine}
    \item \code{/bds/proj1/processed/quality\_report}
  \end{itemize}
\end{frame}

\begin{frame}{Arhitektura}
  3-slojni model:
  \begin{enumerate}
    \item HDFS sloj: skladište (\code{namenode}, \code{datanode})
    \item Spark sloj: obrada (\code{spark-master}, \code{spark-worker-*})
    \item App sloj: submission i CLI (\code{app})
  \end{enumerate}

  \vspace{0.35cm}
  \begin{block}{Tok podataka}
    Raw CSV $\rightarrow$ HDFS \code{/raw} $\rightarrow$ preprocess job $\rightarrow$ HDFS \code{/processed/*}
  \end{block}

  \begin{itemize}
    \item analytics CLI koristi \code{count}, \code{show}, \code{stats}
    \item benchmark kampanje koriste isti HDFS izvor
  \end{itemize}
\end{frame}

\begin{frame}{Docker Compose}
  Servisi:
  \begin{itemize}
    \item \code{namenode}, \code{datanode}
    \item \code{spark-master}, \code{spark-worker-1}, \code{spark-worker-2}
    \item \code{app}
  \end{itemize}

  \vspace{0.35cm}
  Princip:
  \begin{itemize}
    \item isti HDFS izvor za oba benchmark moda
    \item \code{app} container submituje iste komande za \code{local[*]} i \code{spark://...}
  \end{itemize}
\end{frame}

\begin{frame}[fragile]{Upload na HDFS}
  Reproducibilni checkpoint:
  \begin{verbatim}
automation/phase1_hdfs_upload.sh
  \end{verbatim}

  Ključne komande:
  \scriptsize
  \begin{verbatim}
docker compose exec namenode hdfs dfs -mkdir -p /bds/proj1/raw
docker compose exec namenode hdfs dfs -put -f \
  /input/.../jul2018.csv \
  /bds/proj1/raw/jul2018.csv
docker compose exec namenode hdfs dfs -put -f \
  /input/.../dec2018.csv \
  /bds/proj1/raw/dec2018.csv
docker compose exec namenode hdfs dfs -count -h /bds/proj1/raw
  \end{verbatim}
  \normalsize
\end{frame}

\begin{frame}[fragile]{Preprocessing (Spark job)}
  Reproducibilni checkpoint:
  \begin{verbatim}
automation/phase2_preprocessing.sh
  \end{verbatim}

  Job:
  \begin{verbatim}
docker compose exec app /spark/bin/spark-submit \
  --master spark://spark-master:7077 \
  /workspace/scripts/preprocess_ais_hdfs.py \
  --input 'hdfs://namenode:9000/bds/proj1/raw/*.csv' \
  --output-base hdfs://namenode:9000/bds/proj1/processed
  \end{verbatim}
\end{frame}

\begin{frame}[fragile]{Preprocessing logika}
  Ključne transformacije:
  \begin{verbatim}
missing_required = timestamp/vessel_id/lon/lat missing
invalid_geo      = lon/lat out of range
invalid_speed    = speed < 0
invalid_heading  = heading not in [0, 360]
invalid_course   = course not in [0, 360]

clean       <- no hard_fail
quarantine  <- hard_fail rows

clean:       parquet partitioned by year, month
quarantine:  parquet partitioned by reason
quality_report: summary + reason_counts
  \end{verbatim}
\end{frame}

\begin{frame}[fragile]{Analytics CLI}
  Komande:
  \begin{itemize}
    \item \code{count}
    \item \code{show}
    \item \code{stats} (min/max/avg/stddev + count)
  \end{itemize}

  Primeri:
  \begin{verbatim}
ais_analytics.py count --year 2018 --month 12

ais_analytics.py show --year 2018 --month 7 --min-speed 15 \
  --sort-by timestamp --sort-desc --limit 20

ais_analytics.py stats --year 2018 --group-by month --metrics speed
  \end{verbatim}
\end{frame}

\begin{frame}{Benchmark metodologija}
  Pravila za fer poređenje:
  \begin{itemize}
    \item isti dataset: \code{/processed/clean}
    \item ista komanda
    \item režimi: \code{--master local[*]} i \code{--master spark://spark-master:7077}
  \end{itemize}

  \vspace{0.35cm}
  Kampanja:
  \begin{itemize}
    \item 5 iteracija
    \item svaka iteracija: 1 run po modu + poseban report
    \item zatim finalni agregat (5x2)
  \end{itemize}
\end{frame}

\begin{frame}[fragile]{Benchmark alati i artefakti}
  Runner:
  \begin{verbatim}
python3 scripts/benchmark_ais_analytics.py \
  --benchmark-name <name> \
  --modes both \
  --repeats 1 \
  --query-override "count ... | show ... | stats ..."
  \end{verbatim}

  Aggregator:
  \begin{verbatim}
python3 scripts/aggregate_benchmark_runs.py \
  --runs-dir runs/campaign_<id>/runs \
  --output-root runs/campaign_<id>/reports
  \end{verbatim}
\end{frame}

\begin{frame}{Full checkpoint}
  Finalni reproducibilni sistem:
  \begin{center}
    \code{automation/phase5\_full\_benchmarks.sh}
  \end{center}

  \vspace{0.35cm}
  Šta radi:
  \begin{enumerate}
    \item podiže storage + compute + app sloj
    \item puni HDFS i izvršava preprocessing
    \item pokreće 4 benchmark kampanje
    \item za svaku kampanju pravi 5 iteracionih reporta + 1 finalni agregat
  \end{enumerate}
\end{frame}

\begin{frame}[fragile]{Struktura rezultata}
  \scriptsize
  \begin{verbatim}
runs/campaign_<id>/
  runs/
    standalone/<timestamp>/
      meta.json
      results.json
      summary.json
      run_report.md
    cluster/<timestamp>/
      ...
  reports/
    <timestamp>/aggregate.md
    <timestamp>/aggregate.csv
    <timestamp>/aggregate.json
  \end{verbatim}
  \normalsize

  \code{meta.json} čuva raw i processed size snapshot.
\end{frame}

\begin{frame}[fragile]{Rezultati: Stats (speed)}
  Kampanja: \code{campaign\_20260226T101745Z\_stats\_speed}

  \scriptsize
  \begin{verbatim}
stats --year 2018 --group-by month --metrics speed \
  --order-by month --limit 12
  \end{verbatim}
  \normalsize

  \begin{tabular}{lrrrrr}
    \toprule
    mode & n & mean s & stddev s & min s & max s \\
    \midrule
    standalone & 5 & 11.049 & 0.576 & 10.350 & 11.781 \\
    cluster & 5 & 15.444 & 0.458 & 15.023 & 16.202 \\
    \bottomrule
  \end{tabular}

  \vspace{0.25cm}
  Raw: \code{971.54 MiB}, Processed: \code{191.60 MiB}
\end{frame}

\begin{frame}[fragile]{Rezultati: Stats (heading, course)}
  Kampanja: \code{campaign\_20260226T102011Z\_stats\_heading\_course}

  \scriptsize
  \begin{verbatim}
stats --year 2018 --group-by month --metrics heading,course \
  --order-by month --limit 12
  \end{verbatim}
  \normalsize

  \begin{tabular}{lrrrrr}
    \toprule
    mode & n & mean s & stddev s & min s & max s \\
    \midrule
    standalone & 5 & 11.132 & 0.498 & 10.760 & 11.825 \\
    cluster & 5 & 16.069 & 0.611 & 15.484 & 16.734 \\
    \bottomrule
  \end{tabular}

  \vspace{0.25cm}
  Raw: \code{971.54 MiB}, Processed: \code{191.60 MiB}
\end{frame}

\begin{frame}[fragile]{Rezultati: Count}
  Kampanja: \code{campaign\_20260226T102241Z\_count}

  \begin{verbatim}
count --year 2018 --month 12
  \end{verbatim}

  \begin{tabular}{lrrrrr}
    \toprule
    mode & n & mean s & stddev s & min s & max s \\
    \midrule
    standalone & 5 & 9.486 & 0.018 & 9.465 & 9.508 \\
    cluster & 5 & 14.065 & 0.231 & 13.657 & 14.188 \\
    \bottomrule
  \end{tabular}
\end{frame}

\begin{frame}[fragile]{Rezultati: Show}
  Kampanja: \code{campaign\_20260226T102453Z\_show}

  \scriptsize
  \begin{verbatim}
show --year 2018 --month 7 --min-speed 15 \
  --sort-by timestamp --sort-desc --limit 20
  \end{verbatim}
  \normalsize

  \begin{tabular}{lrrrrr}
    \toprule
    mode & n & mean s & stddev s & min s & max s \\
    \midrule
    standalone & 5 & 10.490 & 0.031 & 10.445 & 10.513 \\
    cluster & 5 & 14.309 & 0.220 & 14.167 & 14.697 \\
    \bottomrule
  \end{tabular}
\end{frame}

\begin{frame}{Poređenje cluster vs standalone}
  Cluster/standalone odnos (mean):

  \vspace{0.25cm}
  \begin{tabular}{lrr}
    \toprule
    benchmark & ratio & delta (s) \\
    \midrule
    stats speed & 1.397x & +4.395 \\
    stats heading, course & 1.443x & +4.937 \\
    count dec2018 & 1.483x & +4.579 \\
    show fast jul2018 & 1.364x & +3.819 \\
    \bottomrule
  \end{tabular}
\end{frame}

\begin{frame}{Analiza rezultata}
  Zaključci za trenutni obim (\textasciitilde 1GB raw, \textasciitilde 192MB parquet):
  \begin{itemize}
    \item \code{local[*]} je brži u svim testovima.
    \item Cluster overhead dominira nad dobitkom paralelizacije na ovom obimu podataka.
    \item Stabilnost merenja zavisi od upita: \code{count} i \code{show} su stabilniji, dok \code{stats} ima veću varijansu.
  \end{itemize}

  \vspace{0.35cm}
  Implikacija:
  \begin{itemize}
    \item Za dokaz da cluster postaje bolji, potrebno je povećati ulaz i/ili težinu transformacija.
  \end{itemize}
\end{frame}

\begin{frame}{Ograničenja i sledeći koraci}
  Ograničenja trenutnog eksperimenta:
  \begin{itemize}
    \item jedan obim podataka (\textasciitilde 1GB raw)
    \item jedan host i ograničeni worker resursi
    \item fokus na latency, ne na throughput pod konkurentnim opterećenjem
  \end{itemize}

  \vspace{0.35cm}
  Sledeće:
  \begin{enumerate}
    \item benchmark na većim obimima podataka (npr. 5GB+ raw)
    \item podešavanje Spark parametara
    \item finalna vizuelna analiza pomoću \code{aggregate.csv}
  \end{enumerate}
\end{frame}

\begin{frame}[fragile]{Reproducibilnost}
  Reference komande:
  \begin{verbatim}
# HDFS + preprocess baseline
bash automation/phase2_preprocessing.sh

# jedna benchmark kampanja
python3 scripts/benchmark_ais_analytics.py ...
python3 scripts/aggregate_benchmark_runs.py ...
  \end{verbatim}

  Sve kampanje i izveštaji su timestamped i append-only.
\end{frame}

\begin{frame}
  \centering
  {\Huge Pitanja?}

  \vspace{0.8cm}
  {\Large Kraj}
\end{frame}

\end{document}
